\chapter{The product, concretely}
\label{sec:product}

Everything before this page weighed evidence; this part spends it. The
reader changes here too: the chapters so far served whoever must judge
our verdicts, and the chapters from here serve whoever must build from
them, starting the week the plan is approved. The two jobs meet in one
requirement, so we state it before any design.

\emph{Concrete} has a checkable meaning in this part, and it is not
``more detail.'' A verdict from the decision boxes counts as delivered
only when four anchors exist for it: a contract (which interface,
schema, or component carries it), a test (which entry of the
validation suite would catch its failure), a schedule slot (which tier
and week builds it), and a price (whose weeks pay for it). The master
table of Chapter~\ref{sec:roadmap} already lists every verdict with
its tier and cost; this part's job is to supply the two missing
anchors, contracts and tests, and to close the loop so that no
verdict lives only in a box.

The part runs in build order. This chapter says what the product is
and walks one study through it on day one.
Chapter~\ref{sec:architecture} holds the architectural design, two
interfaces, two adapters, and a database, unchanged from the verdict
it argued. The chapters after it pin the contracts and the validation
suite, and Chapter~\ref{sec:roadmap} closes the part with the
schedule and the bill.

\section{What we ship, and to whom}

The product is an internal Python package: a tuning layer that our
four workload families call, running on Ray Tune as its execution
layer, exactly as Chapter~\ref{sec:architecture} argues. It is not a
service, not a platform, and not a research program; every research
bet it contains (the learning-curve scheduler, the population
flagship) is scheduled behind a working default, not in front of it.

The customers are the projects of Chapter~\ref{sec:workloads}, and
the package meets each one with a configured study rather than a menu:
the routine-budget RL study, the large-parallelism RL study, the
sampler study with its correctness vetoes, and the two-objective
Hamiltonian study. A researcher who knows only their own workload
should be able to run the matching study without reading this part
of the document; the chapters exist so that the person
\emph{extending} the package can reconstruct why each default is what
it is.

What ships, in one sentence per layer: a search-space schema that
knows the three knob kinds, conditional structure, and an optional
prior; a searcher interface and a scheduler interface with wrapped
implementations behind them (Optuna's TPE and the GP searcher on one
side, ASHA and the median rule on the other); a run store that
records configurations, seeds, curves, costs, and lineage; the seed
protocol as the default study shape rather than an option; and the
validation suite, which is a deliverable of the same rank as the
tuner, because every falsification promise made in
Chapters~\ref{sec:modelfree} through~\ref{sec:transfer} is an entry
in it.

\section{A study on day one}

The fastest way to make the product concrete is to run one study
through it in words, at tier~0, using only components that exist and
are maintained today plus our thin glue. Nothing in this walk waits
on research; that is what ``wire up what exists'' meant in
Chapter~\ref{sec:roadmap}, and the walk doubles as the acceptance
test for tier~0.

A researcher opens a study for the running example: PPO on the Brax
task, routine budget, the situation of Chapter~\ref{sec:workloads}.
They declare the fifteen-knob space of Chapter~\ref{sec:problem} in the
schema, unchanged and with no knob added: learning rate log-uniform
between $10^{-5}$ and $10^{-2}$, the ordinal widths, depths, and batch
sizes as ordered integers, and the two spectral-normalization flags as
categories. The BG-PBT space carries no conditional knob, so this study
declares a flat space; the schema's conditional support is exercised by
the annealing switch of Chapter~\ref{sec:problem} and by the off-policy
and sampler workloads of Chapter~\ref{sec:workloads}, and the
capability test for it lives there rather than here. The schema also
accepts the folklore prior (learning rate
near $3 \times 10^{-4}$, discount near $0.99$) as a first-class
field. On day one nothing consumes it; recording it from day one is
deliberate, because Chapter~\ref{sec:transfer} showed that priors
are cheap to store, expensive to retrofit, and eventually worth
real trials.

They declare the objective explicitly, mean evaluated return aggregated
over the declared tuning seeds, because Chapter~\ref{sec:workloads}
showed that final return, area under the curve, and truncated return
genuinely disagree, and a hidden default here silently changes every
ranking downstream. This tuning-seed mean is what the searcher observes
and optimizes. After selection is frozen, the winner (or Pareto front)
is evaluated once on disjoint test seeds that the tuner never saw, and
that held-out performance is what the final report quotes.

This seed protocol is the study's default shape: search driven by
tuning-seed feedback, final evaluation on test seeds. This is not an
option the researcher enables; it is the study shape, and opting
\emph{out} is the act that requires writing something down. The
evidence for that inversion is the seed-noise record of
Chapters~\ref{sec:problem} and~\ref{sec:workloads}, and the protocol
is what makes the test set a genuine held-out evaluation rather than
part of the search feedback loop.

The searcher on day one is the wrapped TPE; the scheduler is ASHA
with the median rule beside it, granting first rungs at one tenth of
full fidelity. While the study runs, every trial streams its curve
into the run store with its configuration, seed, rung history, and
bill, and the store computes one diagnostic the researcher did not
ask for: the rank correlation between rung scores and final scores,
logged per study. That number is Chapter~\ref{sec:multifidelity}'s
crossing-curves hazard turned into a measured quantity of our own
workloads, and the validation suite consumes it later.

What the researcher gets back is the incumbent configuration with
its test-seed mean and spread, the full record of trials in the
store, and the study's own audit trail: which trials were stopped at
which rungs, and what the early-versus-final correlation was. What
they do not get on day one, stated as plainly: no schedules (the
population line is tier~2), no premium multi-objective searcher (the
Hamiltonian study runs on the wrapped evolutionary default until the
qLogNEHVI searcher lands), and no prior-weighted acquisition yet.
Each absence has a tier and a week in Chapter~\ref{sec:roadmap};
none of them blocks the day-one study from being useful, which is
the point of the tier ordering.

\section{The register, and how this part uses it}

Chapter~\ref{sec:roadmap}'s Table~\ref{tab:decisions} is the spine
of this part: every verdict, its tier, its cost, and the chapter
holding its evidence. The two chapters that follow extend each row
with its remaining anchors. The contracts chapter names the
interface or schema field that carries the verdict, so that a reader
can point from ``priors are first-class'' to the line of the schema
that makes it true. The validation chapter names the suite entry
that would catch the verdict failing, and it inherits its content
from promises already made: every ``argued from four angles''
section in Chapters~\ref{sec:modelfree}
through~\ref{sec:transfer} ended with a falsification criterion, and
the suite is those criteria collected, priced, and scheduled rather
than restated.

Two bookkeeping rules keep the register honest while this part is
written. A verdict with no plausible suite entry is marked
untestable in place, with a sentence saying why, rather than
receiving a decorative test. And any design sentence in this part
that cannot cite a chapter of evidence is labeled an engineering
judgment, so the reader always knows which statements the survey
backs and which ones the build team owns.

\section{What would overturn this chapter}

The day-one walk is this chapter's load-bearing claim: tier~0 is
composition, not construction. It fails, and this part gets
rewritten, if the walk cannot be reproduced in software within the
tier~0 budget of Chapter~\ref{sec:roadmap}, if any step turns out to
require a component the audit of Chapter~\ref{sec:libraries} did not
verify as maintained, or if the seed-protocol default proves too
expensive to be the default on our routine studies. The first
validation-suite run is scheduled to test exactly this walk, and the
suite, not this chapter, has the last word.
